\documentclass[11pt,a4paper]{article}
\usepackage[margin=2.2cm]{geometry}
\usepackage{amsmath,amssymb}
\usepackage{graphicx}
\usepackage{booktabs}
\usepackage[font=small,labelfont=bf]{caption}
\usepackage{xcolor}
\usepackage[colorlinks=true,linkcolor=blue!50!black,urlcolor=blue!50!black]{hyperref}
\usepackage{enumitem}
\setlist{itemsep=2pt,topsep=4pt}

\graphicspath{{../exp_data/}}
\newcommand{\fig}[1]{\includegraphics[width=\textwidth]{#1}}

\title{A house cracking under wind: phase-field fracture with the corrected Green preconditioner}
\author{muFFTTO, \texttt{experiments/linear\_elastic\_mesh\_study}}
\date{24 October 2026}

\begin{document}
\maketitle

\noindent\fbox{\parbox{\dimexpr\textwidth-2\fboxsep-2\fboxrule}{%
\textbf{Summary.}
\begin{itemize}
  \item A masonry house (6 $\times$ 8\,m, ridge 5\,m, walls 0.25\,m) on voxels, clamped at the foundation, under
        self-weight and growing wind, with AT2 phase-field fracture. Run on the GPU, controlled by the eaves displacement
        so that the softening branch after the peak wind is followed.
  \item \textbf{Where it cracks.} With self-weight: diagonal shear cracks in both gable walls, from the windward
        eaves corner down to the gable windows (2--3\,m height), later a band near the window sills (0.5--1\,m).
        Without self-weight the cracks start low: at 0.5--1.5\,m in the gable walls and at the windward corners.
        Peak wind 10.0\,kPa (with self-weight) and 8.8\,kPa (without).
  \item \textbf{A bug in the Green preconditioner} made every 3D solve stall in physical units ($E = 10$\,GPa): the
        un-inverted zero-frequency block was left at its round-off. Fixed in \texttt{domain.py}, with a regression test
        (Section~\ref{sec:bug}).
  \item \textbf{Which preconditioner.} On the same matrices (frozen damage, intact to cracked) the corrected
        Green needs 245--370 CG iterations, plain Green 360--500, Green-Jacobi 600--890, and a dual with the corrected
        Green and the void taken as a soft material inside the dual 124--187. A dual iteration costs twice as much, so
        corrected Green and this dual are equally fast, 1.2--1.4 times faster than plain Green and 2.0--2.3 times faster
        than Green-Jacobi. The dual with the void as zero stiffness does not converge at all (Section~\ref{sec:precond}).
\end{itemize}}}

%-----------------------------------------------------------------------------------------------------------------------
\section{Set-up}
\label{sec:setup}

\paragraph{Geometry.} The house is 6\,m deep in the wind direction $x$, 8\,m long in $z$, with the eaves at 3\,m and
the ridge at 5\,m (Figure~\ref{fig:geometry}). Walls and roof slab are 0.25\,m thick masonry ($E = 10$\,GPa,
$\nu = 0.2$, 1800\,kg/m$^3$); the roof overhangs by 0.25\,m. The door and the seven windows are closed by soft panels
($E = 0.2$\,GPa, 500\,kg/m$^3$) that pass their wind load to the frames. Around the house is a void of 0.5\,m (zero
stiffness). The walls are clamped at the foundation: every node of the bottom face of a wall voxel.

\paragraph{Loads.} Self-weight at full value, and the wind with Eurocode-style external pressure coefficients on the
outer faces, $p = c_{pe}\,q$: windward wall $+0.8$ (pushed in), leeward wall $-0.5$, side walls $-0.7$ (sucked out), roof
$-0.6$ (uplift). A voxel face counts as outer if the straight line from it out of the house meets no solid. At
$q = 1$\,kPa the wind gives 56\,kN horizontally and 33\,kN uplift; the self-weight is 657\,kN.

\begin{figure}[h]
  \fig{house_geometry_h0.125.png}
  \caption{The house at voxel size 12.5\,cm (grid $60 \times 51 \times 76$), with a sample of the wind loads: red
  windward pressure, blue leeward suction, purple side suction, green roof uplift.}
  \label{fig:geometry}
\end{figure}

\paragraph{Fracture.} AT2 phase field as in \texttt{experiments/dual\_preconditioner/damage\_AT2\_small\_strain}:
$\mathbb{C}(d) = ((1-d)^2 + 10^{-4})\,\mathbb{C}_0$ (hybrid formulation), history field
$H = \max_t \psi^+(\varepsilon)$ with the volumetric--deviatoric split of Amor et al., staggered scheme with
tolerance $\max|d - d_{\mathrm{old}}| < 10^{-2}$. Only the masonry can crack. The length scale is two voxels,
$\ell = 2h$, and the toughness follows from a tensile strength $f_t = 0.2$\,MPa through the AT2 peak stress
$\sigma_c = \tfrac{9}{16}\sqrt{E G_c/(3\ell)}$: $G_c = 9.5$\,N/m at $h = 12.5$\,cm.

\paragraph{Load control.} Under growing wind (force control) the house breaks in one load step once the peak is passed:
there is no equilibrium beyond it. The run is therefore controlled by the mean horizontal displacement of the eaves.
For a fixed damage field the problem is linear, $u = u_g + \lambda u_w$ (self-weight, wind of 1\,kPa), so in every
staggered iteration both parts are solved (warm started) and the wind factor $\lambda$ is chosen to hit the prescribed
eaves displacement. The prescribed displacement grows in steps adapted to the number of staggered iterations. The scale
is $q_{\mathrm{crit}}$, the wind at which the largest principal stress of the linear solution reaches $f_t$.

\paragraph{Solvers.} CG with energy-norm stopping (relative $10^{-6}$, delay 10) for both problems. Elasticity: the
corrected Green preconditioner (exact inverse of the homogeneous reference problem with the clamped dofs, $\mathbf{S}$
assembled once from 1944 Green applications in 10\,s; it never changes, since the damage changes only the material).
Damage: Green-Jacobi.

\paragraph{Output.} \texttt{exp\_data/house/damage\_h0.125\_ft0.2.nc} (NetCDF): one frame per staggered iteration with
the damage and displacement fields, the wind, the eaves displacement, maximum damage, cracked masonry volume, CG
iterations and times, and the full CG convergence histories (energy-error estimate per iteration) of both solves.

%-----------------------------------------------------------------------------------------------------------------------
\section{Results}
\label{sec:results}
Two runs at voxel size 12.5\,cm ($60 \times 51 \times 76$ voxels, 0.7 million unknowns), one with and one without
self-weight (Figure~\ref{fig:curves}). Both start elastic, soften as the damage grows, and break suddenly at the peak
wind: the first crack runs through in one displacement step (95 and 68 staggered iterations), the wind drops to 60\,\%
and 42\,\% of the peak, and the house then softens further as the cracks extend.

\begin{figure}[h]
  \fig{house_damage_curves_h0.125.png}
  \caption{Left: wind against the controlled eaves displacement; the thin lines are the staggered iterations inside the
  snaps. Middle: cracked masonry (voxels with mean nodal damage above 0.9). Right: CG iterations per elastic solve over
  all staggered iterations.}
  \label{fig:curves}
\end{figure}

\begin{center}
\small
\begin{tabular}{lcc}
  \toprule
  & with self-weight & without self-weight \\
  \midrule
  linear first crack $q_{\mathrm{crit}}$ (max principal stress $= f_t$) & 2.22\,kPa at 2.3\,m height & 1.84\,kPa at 2.2\,m \\
  peak wind & 10.0\,kPa & 8.8\,kPa \\
  first crack: wind after the snap, cracked masonry & 6.0\,kPa, 0.20\,m$^3$ & 3.7\,kPa, 0.25\,m$^3$ \\
  last converged step: wind, cracked masonry & 5.7\,kPa, 0.29\,m$^3$ & %NOGRAV_LAST \\
  where & gables 2--3\,m (diagonal), 0.5--1\,m & gables and windward corners, 0.5--1.5\,m \\
  wall time (GPU, shared from 03:00) & 5\,h (stopped at the deadline) & %NOGRAV_TIME \\
  \bottomrule
\end{tabular}
\end{center}

\paragraph{With self-weight} the base is precompressed (about 60\,kPa at the foundation), so the first crack forms
where the wind stresses are highest relative to that: at the re-entrant corners near the eaves of the two gable walls,
which carry the wind as shear walls. It grows as a diagonal crack from the windward eaves corner down to the gable
window (Figure~\ref{fig:last}). Later a second band opens near the window sills. In the last load step (85) the
remaining capacity collapses: the wind falls to zero within the staggered iterations, which did not converge before
the deadline (06:00) stopped the run; the NetCDF holds all of them.

\paragraph{Without self-weight} uplift and overturning put the lower walls in tension, and the cracks form low: in the
gable walls at 0.5--1.5\,m and at the windward corners where the facade meets them.

\paragraph{Resolution.} At 25\,cm ($\ell = 0.5$\,m, two voxels through the wall thickness counted as one) the same
house with self-weight cracked in a horizontal band at the window sills (0.75--1\,m), peak wind 15.3\,kPa: the coarse
crack band smears out the corner concentrations. At 12.5\,cm the corners near the eaves win and the peak is lower.
The crack pattern is therefore not yet converged in $h$; 6.25\,cm (5.6 million unknowns) would be the next check.

\paragraph{Videos.} \texttt{exp\_data/house\_damage\_h0.125\_ft0.2.mp4} and
\texttt{exp\_data/house\_damage\_h0.125\_ft0.2\_nograv.mp4}: one frame per converged step, plus every staggered
iteration of the steps that needed more than five, so the cracks are seen running through during the snaps.

\begin{figure}[h]
  \fig{house_damage_h0.125_ft0.2_last.png}
  \caption{With self-weight, the last frame: deformed house (magnified) from the windward and the leeward side,
  masonry coloured by the damage; wind against eaves displacement; CG iterations.}
  \label{fig:last}
\end{figure}

\begin{figure}[h]
  \fig{house_damage_h0.125_ft0.2_nograv_last.png}
  \caption{As Figure~\ref{fig:last}, without self-weight.}
  \label{fig:last_nograv}
\end{figure}

%-----------------------------------------------------------------------------------------------------------------------
\section{Which preconditioner for the cracking house?}
\label{sec:precond}
The elastic solves dominate the cost. To compare preconditioners fairly, the damage field of five converged steps of
the run with self-weight (intact, $\max d = 0.2$, peak, after the first crack, last) was frozen and the same elastic
problem solved from zero with each preconditioner (\texttt{house\_precond\_study.py}, Figure~\ref{fig:precond}). All
give the same compliance to seven digits.

\begin{figure}[h]
  \fig{house_precond_h0.125.png}
  \caption{CG iterations, time per solve and time per iteration, five frozen damage states, 12.5\,cm. Hatched:
  projected (Dirichlet by $\mathbf{Q}$); solid: corrected for the clamped foundation. Times taken while the second
  run shared the GPU.}
  \label{fig:precond}
\end{figure}

\begin{center}
\small
\begin{tabular}{lccccc}
  \toprule
  CG iterations & intact & $\max d$ 0.2 & peak & first crack & last \\
  \midrule
  Green, projected                  & 375 & 369 & 358 & 445 & 497 \\
  Green, corrected                  & 248 & 245 & 257 & 315 & 368 \\
  Green-Jacobi, projected           & 615 & 601 & 637 & 754 & 825 \\
  Green-Jacobi, corrected           & 677 & 665 & 678 & 818 & 890 \\
  dual, corrected, void as soft     & \textbf{128} & \textbf{124} & \textbf{133} & \textbf{161} & \textbf{187} \\
  \midrule
  time [s]: Green, corrected        & 27.5 & 27.0 & 29.2 & 35.9 & 41.7 \\
  time [s]: dual, void as soft      & 27.9 & 27.0 & 29.0 & 35.1 & 40.7 \\
  time [s]: Green, projected        & 38.6 & 37.9 & 36.8 & 45.8 & 50.8 \\
  time [s]: Green-Jacobi, projected & 63.1 & 61.9 & 65.6 & 77.7 & 85.1 \\
  \bottomrule
\end{tabular}
\end{center}

\begin{itemize}
  \item \textbf{Green beats Green-Jacobi here.} The structure has only two stiff phases (masonry, panels) and void,
        all voxel-aligned, so the Green reference fits well; the Jacobi scaling does not help: Green-Jacobi needs 1.6--1.8 times the
        iterations of plain Green and 2.2--2.5 times those of the corrected Green. Damage hardly changes this: all methods need 30--50\,\% more iterations once cracks are open.
  \item \textbf{The correction for the clamped foundation} saves a third of the iterations of Green (as in the 2D
        study), but makes Green-Jacobi slightly worse (by 5--10\,\%); the correction is exact for the
        $2\mu\,\mathbb{I}$ reference, not for the Jacobi-scaled operator that CG actually sees.
  \item \textbf{The dual} $\mathbf{M}_c\,\mathbf{K}(\mathbb{C}_{\mathrm{ref}}\,\mathbb{C}_d^{-1}\,
        \mathbb{C}_{\mathrm{ref}})\,\mathbf{M}_c$ halves the iterations again but needs two Green and one stiffness
        application per iteration, so it is as fast as the corrected Green. It needs $\mathbb{C}_d^{-1}$ in the void:
        with the pseudo-inverse (zero in the void) it reduces to $\mathbf{G}\mathbf{K}_{\mathrm{solid}}\mathbf{G}$,
        which squares the spectrum instead of inverting it, and CG did not converge in 6000 iterations (25\,cm, both
        projected and corrected). Taking the void as $\varepsilon\,\mathbb{C}_{\mathrm{ref}}$ inside the dual only (the
        problem keeps zero stiffness) fixes this, and the result hardly depends on $\varepsilon$: 149--220 iterations at
        25\,cm for $\varepsilon = 10^{-2}, 10^{-3}, 10^{-4}$.
  \item \textbf{Damage solves} (Green-Jacobi) need 13--30 CG iterations and under a second; they are not the cost.
\end{itemize}

%-----------------------------------------------------------------------------------------------------------------------
\section{The bug: Green in physical units}
\label{sec:bug}

The first 3D solves stalled: the energy error fell only like $1/k$, plain and corrected Green converged to different
solutions, and the Ritz values were impossible ($\lambda_{\max} \approx 3\times10^5$ for the corrected Green, negative
ones for plain Green). The stiffness operator and the raw Green operator applied with $E = 1$ were exactly symmetric;
with $E = 10$\,GPa the Green operator was not: $x \cdot Gy \ne y \cdot Gx$ and $x \cdot Gx < 0$ for random fields.

\paragraph{Cause.} \texttt{get\_preconditioner\_Green\_mugrid} inverts the blocks $\hat K(k)$ of the reference stiffness
for every frequency except $k = 0$, which is singular; that block was left as it was, assuming it is zero. It is zero
only up to round-off, about $|\mathbb{C}_{\mathrm{ref}}|\,\varepsilon \approx 10^{10}\cdot 10^{-16}\cdot N$, while the
inverted blocks are of order $1/|\mathbb{C}_{\mathrm{ref}}| \approx 10^{-10}$. In physical units the round-off dominated
the operator. With $E = 1$, as in all earlier studies, it is harmless, which is why it never showed.

\paragraph{Fix.} The zero-frequency block is set to exactly zero (one line in \texttt{muFFTTO/domain.py}, the one-node
branch; the multi-node branch already used a pseudo-inverse with a relative cut-off). The new test
\texttt{test\_green\_preconditioner\_symmetric\_for\_physical\_scale} in \texttt{test/test\_preconditioners.py} checks
symmetry, positivity and $G K x = x - \bar x$ for $E = 1$ and $10^{10}$ in 2D and 3D; it fails without the fix.
After the fix: $\lambda_{\max} = 0.968$ for the corrected Green, 1.28 for plain Green.

\paragraph{Green-Jacobi} needs a dimensionless reference ($2\mu/E\,\mathbb{I}$): the Jacobi scaling already carries the
material scale.

%-----------------------------------------------------------------------------------------------------------------------
\section*{Files}
All in \texttt{experiments/linear\_elastic\_mesh\_study/}.
{\raggedright
\begin{description}[font=\ttfamily\small,style=nextline]
  \item[house\_wind\_3D.py] geometry, wind loads, the elastic solve and the geometry figure (\texttt{GEOMETRY=1}).
  \item[house\_damage\_3D.py] the fracture run (GPU), NetCDF output.
  \item[house\_damage\_video.py] the video and its last frame from the NetCDF.
  \item[house\_precond\_study.py] the frozen-damage preconditioner comparison and its figure (\texttt{PLOT=1}).
  \item[house\_elastic\_diag.py] diagnostics of the elastic solve (symmetry, power iteration, variants).
  \item[house\_geometry\_video.py] rotating view of the geometry.
\end{description}
}
\end{document}
